\bookchapter{الفصل 41: المشهد الكامل لأدوات البرمجة بالذكاء الاصطناعي في 2026}

\chapterintro{لنبدأ بتحذير وقائي: هذا الفصل «لقطة» من منظومة الأدوات كما كانت في أغسطس 2026، وستصبح قديمة بعد نصف عام — وهذا بالضبط هو الدرس الأول الذي يريد هذا الفصل أن يعلّمه لك. هدف الفصل ليس أن يعلّمك «أي أداة تستخدم»، بل أن يساعدك على بناء إطار حكم لـ «كيف تختار الأدوات». وحين تفهم الفلسفة التصميمية للأدوات المختلفة، ستستطيع في أي وقت اتخاذ الخيار المناسب لك.}

\booksection{أولًا: المدارس الثلاث الكبرى وإطار الاختيار}

تنقسم أدوات البرمجة بالذكاء الاصطناعي في 2026 تقريبًا إلى ثلاث مدارس كبرى، وخلف كل مدرسة فلسفة تصميمية مختلفة. إليك اللوحة الشاملة أولًا:

\begin{center}
\begin{tabularx}{\linewidth}{|X|X|X|X|}
\hline
\textbf{المدرسة} & \textbf{الأدوات الممثلة} & \textbf{الفلسفة التصميمية} & \textbf{الحالات المناسبة} \tabularnewline
\hline
ذكاء اصطناعي مدمج في IDE & Cursor 3, Trae, Windsurf, CodeBuddy & دمج قدرات الذكاء الاصطناعي بعمق داخل المحرر & التطوير اليومي، ومراجعة الشيفرة، والتحسين التدريجي \tabularnewline
وكيل الطرفية (Agent) & Claude Code, Codex CLI, Aider & يتخذ الطرفية واجهةً له، فيعالج الذكاء الاصطناعي نظام الملفات والأوامر مباشرة & العمليات على مستوى المشروع، والتعديلات الجماعية، والتكامل مع CI/CD \tabularnewline
باني سحابي (Builder) & v0, Bolt, Replit Agent & إنشاء مشاريع كاملة من الصفر داخل المتصفح & النماذج الأولية السريعة، والبدء دون إعدادات، والمعاينة الفورية \tabularnewline
\hline
\end{tabularx}
\end{center}

\begin{center}\small\color{gray}المدارس الثلاث الكبرى لأدوات البرمجة بالذكاء الاصطناعي\end{center}

مفتاح اختيار الأداة ليس «أيها الأفضل»، بل «أيها الأنسب لحالتك الحالية». وإليك ثلاثة معايير عملية للحكم:

\begin{itemize}
  \item ما مستواك في البرمجة؟ إن كنت مبتدئًا تمامًا فاختر الباني السحابي، وإن كانت لديك أساسيات فاختر الذكاء الاصطناعي المدمج في IDE، أما المستخدمون المتقدمون فيختارون وكيل الطرفية.
  \item في أي مرحلة يقع مشروعك؟ للبدء من الصفر اختر الباني السحابي، وللتطوير التكراري اختر الذكاء الاصطناعي المدمج في IDE، وللصيانة وإعادة الهيكلة اختر وكيل الطرفية.
  \item ما احتياجاتك من النشر؟ للتحقق السريع اختر الباني السحابي، وللمشاريع الجادة اختر الذكاء الاصطناعي المدمج في IDE، ولأتمتة سير العمل اختر وكيل الطرفية.
\end{itemize}

\booksection{ثانيًا: مدرسة الذكاء الاصطناعي المدمج في IDE}

الذكاء الاصطناعي المدمج في IDE هو المدرسة الأكثر انتشارًا: تُزرع قدرات الذكاء الاصطناعي داخل محرر الشيفرة لديك، فيُحتفظ بتجربة البرمجة التقليدية وتُضاف فوقها تحسينات بالذكاء الاصطناعي. وفي 2026 أبرز ممثلي هذه المدرسة هم Cursor 3 وTrae وWindsurf وCodeBuddy.

\booksubsection{Cursor 3}

Cursor هي بيئة التطوير للبرمجة بالذكاء الاصطناعي الأوسع استخدامًا في العالم حاليًا. ويشكل Cursor 3 الصادر في أبريل 2026 نقطة تحول مهمة لأدوات البرمجة بالذكاء الاصطناعي: إذ تجاوز معدل استخدام وضع الوكيل (Agent) لأول مرة وضع الإكمال بمفتاح Tab — فقد انتقل المطورون من «الذكاء الاصطناعي يكمل لي الشيفرة» إلى «الذكاء الاصطناعي ينجز لي المهمة». وتشمل المزايا الأساسية لـ Cursor 3: التحرير متعدد الملفات (يستطيع الذكاء الاصطناعي تعديل عدة ملفات في آن واحد لإنجاز المهمة)، والوعي بالسياق (فهم تلقائي لبنية المشروع وعلاقات الاعتماد)، والتكامل العميق مع الطرفية (يستطيع الذكاء الاصطناعي تشغيل الأوامر وعرض المخرجات مباشرة من داخل المحرر).

\booksubsection{Trae}

Trae هي بيئة تطوير للبرمجة بالذكاء الاصطناعي أطلقتها شركة ByteDance، ولها قاعدة مستخدمين كبيرة في السوق الصيني. ويجسد وضع SOLO لديها (المعروف أيضًا بوضع Builder) قدرة «إنشاء مشروع بجملة واحدة» النموذجية — إذ تصف المتطلبات، ويُنجز الذكاء الاصطناعي كامل المسار من تهيئة المشروع إلى توليد الشيفرة بنفسه. وتكمن ميزة Trae في دعمها للغة الصينية وتلائمها لبيئة الشبكة داخل الصين، ما يجعلها مناسبة للمستخدمين هناك للبدء السريع.

\booksubsection{أدوات أخرى}

تشتهر Windsurf (وكانت تسمى Codeium) بتجربة إكمال شيفرة سلسة، وهي مناسبة للمطورين الباحثين عن تجربة «الكتابة بسرعة الطيران». أما Tencent Cloud CodeBuddy فبوصفه إضافة لـ VS Code يجيد التكامل مع المنظومة السحابية، وهو مناسب للمشاريع التي تخطط للنشر على Tencent Cloud أو لتطوير برامج WeChat المصغرة. وتوفر Tongyi Lingma من Alibaba وComate من Baidu تجربة تكامل مماثلة مع بيئات التطوير، ولكلٍّ منهما خصائص خاصة في تحسين السيناريوهات الصينية.

\booksection{ثالثًا: مدرسة وكيل الطرفية}

وكيل الطرفية هو المدرسة الأسرع نموًا في 2026. وفكرتها الجوهرية: لا ينبغي أن يكون الذكاء الاصطناعي مجرد إضافة لمحرر، بل ينبغي أن يستطيع التعامل مع بيئة التطوير كاملة كما يفعل المطور — قراءة الملفات وكتابتها، وتشغيل الأوامر، وعرض المخرجات، وتثبيت الاعتماديات. وClaude Code هو الممثل الأبرز لهذه المدرسة.

\booksubsection{Claude Code}

Claude Code أداة برمجة بالذكاء الاصطناعي تعمل في الطرفية أطلقتها شركة Anthropic، وقد بلغت إيراداتها السنوية 2.5 مليار دولار بحلول منتصف 2026. وطريقة عملها تختلف كليًا عن بيئات التطوير التقليدية: تصف المهمة بلغة طبيعية في الطرفية، فيقرأ Claude Code الملفات المعنية بنفسه ويعدل الشيفرة ويشغل الاختبارات ويفحص النتائج ويجري التصحيحات تكراريًا — وكل ذلك دون حاجة منك إلى فتح المحرر.

تكمن ميزة Claude Code في «الوعي بالمشروع كاملًا»: فهو غير مقيَّد بنافذة المحرر، ويمكنه التجوال بحرية في مجلد المشروع بأكمله وفهم علاقات الاعتماد بين الملفات. وهو يدعم أيضًا ملف ذاكرة المشروع CLAUDE.md، ليعرف الذكاء الاصطناعي خلفية المشروع منذ بداية كل جلسة. وفي مهام مثل إعادة هيكلة المشاريع الكبيرة والتعديلات الجماعية وكتابة الاختبارات، يكون وكيل الطرفية أكثر كفاءة من إضافات بيئات التطوير.

\booksubsection{Codex CLI وAider}

يوفر Codex CLI من OpenAI تجربة برمجة مشابهة في الطرفية، ويدعم ميزة /import التي تتيح استيراد مكتبة شيفرة قائمة ليفهمها الذكاء الاصطناعي. أما Aider فهي أداة مفتوحة المصدر للبرمجة بالذكاء الاصطناعي في الطرفية، تدعم نماذج متعددة (GPT وClaude وDeepSeek وغيرها)، وتناسب المستخدمين الراغبين في تبديل النماذج بمرونة. ولهذه الأدوات خصائصها المتباينة لكن فكرتها الجوهرية واحدة: تمكين الذكاء الاصطناعي في الطرفية من القدرة الكاملة على التعامل مع بيئة التطوير.

\booksection{رابعًا: مدرسة الباني السحابي}

الباني السحابي هو المدرسة الأنسب للمستخدمين المبتدئين تمامًا. فلا تحتاج إلى تثبيت أي برنامج، ويكفي أن تفتح المتصفح وتصف متطلباتك، فيُنشئ الذكاء الاصطناعي مشروعًا كاملًا في السحابة ويوفر معاينة فورية.

\begin{itemize}
  \item v0 (من إنتاج Vercel): يجيد توليد مكونات واجهة المستخدم والصفحات، ومناسب لصنع نماذج أولية لصفحات الويب بسرعة.
  \item Bolt: يدعم إنشاء تطبيق متكامل من جملة واحدة، يشمل الواجهة الأمامية والخلفية وقاعدة البيانات.
  \item Replit Agent: يعمل داخل بيئة التطوير السحابية من Replit، ولا يقتصر دوره على إنشاء المشاريع بل يستطيع أيضًا تشغيلها ونشرها مباشرة في السحابة.
\end{itemize}

الميزة المشتركة لهذه الأدوات هي «البدء دون أي إعدادات» — تفتح المتصفح وتبدأ فورًا. أما قيودها فهي أن بيئة التحرير السحابية تصبح أقل مرونة من بيئة التطوير المحلية حين يتعقد المشروع. لذلك فهي الأنسب لمرحلة النموذج الأولي «من الصفر إلى الواحد»: استخدم الباني السحابي للتحقق السريع من الفكرة، وبعد نجاح التحقق انقل العمل إلى بيئة تطوير محلية للتطوير المتعمق.

\booksection{خامسًا: الأدوات تتغير والمنهجية تبقى}

ثلاث عادات ينبغي أن تحافظ عليها حتى لو اختفت غدًا كل أدوات البرمجة بالذكاء الاصطناعي: فكّر بوضوح قبل العمل، واعمل خطوة صغيرة واحدة في كل مرة، وتحقق بعد الإنجاز. الأدوات مضخّم والمنهجية أساس. المضخّم يمكن استبداله، أما الأساس فلا يهتز.

لنستعرض تحولات الأدوات من 2024 إلى 2026: في 2024 كان GitHub Copilot هو الملك، وفي 2025 حلّ Cursor محله، وفي 2026 برز Claude Code فجأة بقوة. ومع ظهور كل أداة جديدة يعلن بعضهم أن «هذه المرة الأمر مختلف»، لكن يتبين في النهاية أن الذين يتقنون الأداة الجديدة هم بالضبط أصحاب المنهجيات الراسخة. فهم لا يحتاجون عند تبديل الأداة إلى سوى أيام قليلة من التأقلم، بينما من تعلم تشغيل الأداة فقط يبدأ التعلم من جديد مع كل تبديل.

لذلك فإن مسار التعلم الذي يقترحه هذا الكتاب هو: ابدأ بأداة واحدة (أي أداة كانت) واختبر في الممارسة مبادئ التفكير الخمسة الكبرى؛ وحين تصبح منهجيتك راسخة بما يكفي، جرّب أدوات من مدارس مختلفة واستشعر الفروق في فلسفتها التصميمية؛ وفي النهاية ستكوّن مجموعتك الخاصة من الأدوات — مثل استخدام Cursor للتطوير اليومي، وClaude Code لإعادة الهيكلة والاختبارات، وv0 للنماذج الأولية السريعة. تختلف مجموعة الأدوات من شخص لآخر، لكن المنهجية واحدة مشتركة.

\begin{tcbpractice}{هيّا نطبق: جهّز لنفسك «طقم الأدوات» المناسب}
حدد موقعك بالمعايير الثلاثة (أساسيات البرمجة، ومرحلة المشروع، واحتياجات النشر): إن كنت مبتدئًا تمامًا جرّب Bolt أو v0 لصنع صفحة؛ وإن كانت لديك أساسيات نفّذ مهمة تكرارية باستخدام Cursor أو Trae؛ وإن أردت التحدي شغّل Claude Code مرة واحدة في الطرفية. ثم أجب: أي مدرسة تناسب مرحلتك الحالية أكثر؟

\end{tcbpractice}

\begin{tcbsummary}{خلاصة الفصل}
\begin{itemize}
  \item المدارس الثلاث: الذكاء الاصطناعي المدمج في IDE (التطوير اليومي)، ووكيل الطرفية (العمليات على مستوى المشروع)، والباني السحابي (نماذج أولية للمبتدئين تمامًا)
  \item اختيار الأداة بثلاثة أسئلة: أساسيات البرمجة؟ مرحلة المشروع؟ احتياجات النشر؟
  \item وضع الوكيل في Cursor 3 تفوق لأول مرة على الإكمال بمفتاح Tab؛ وإيرادات Claude Code السنوية بلغت 2.5 مليار دولار
  \item الباني السحابي مناسب لمرحلة «من الصفر إلى الواحد»، وبعد التحقق تنتقل إلى بيئة التطوير المحلية للتطوير المتعمق
  \item تتكرر دورة تحديث الأدوات كل ثلاثة أشهر: من يتعلم المنهجية لا يحتاج عند تبديل الأداة إلى سوى أيام قليلة من التأقلم
\end{itemize}

\end{tcbsummary}
